% !TEX program = lualatex
% !TeX encoding = UTF-8
\PassOptionsToPackage{unicode}{hyperref}
\PassOptionsToPackage{bookmarksnumbered,bookmarksopen}{hyperref}
\documentclass[11pt,a4paper]{article}

% ============================================================
% 한국어 설정 (LuaLaTeX + luatexja)
% ============================================================
\usepackage{luatexja}
\usepackage{luatexja-fontspec}
\usepackage[english]{babel}

% 한국어 고딕체 (sans) – Noto Sans KR
\setsansjfont{Noto Sans KR}[
UprightFont = * Regular,
BoldFont = * Bold,
Script = Hangul,
Language = Korean
]

% 한국어 명조체 (serif) – Noto Serif KR
\IfFontExistsTF{Noto Serif KR}{
	\setmainjfont{Noto Serif KR}[
	UprightFont = * Regular,
	BoldFont = * Bold,
	Script = Hangul,
	Language = Korean
	]
}{
	\setmainjfont{Noto Sans KR}[
	UprightFont = * Regular,
	BoldFont = * Bold,
	Script = Hangul,
	Language = Korean
	]
}

% 고정폭 폰트 – 라틴 문자는 Latin Modern Mono, 한글은 Noto Sans KR
\setmonojfont{Noto Sans KR}[
UprightFont = * Regular,
BoldFont = * Bold,
Script = Hangul,
Language = Korean
]

% 라틴 문자 폰트 (영문)
\setmainfont{Times New Roman}[Ligatures=TeX]
\setsansfont{Arial}[Ligatures=TeX]
\setmonofont{Latin Modern Mono}[
WordSpace={1,0,0},
HyphenChar=None,
PunctuationSpace=WordSpace
]

% ============================================================
% 패키지 (원본과 동일)
% ============================================================
\usepackage{amsmath,amssymb,amsfonts}
\usepackage{geometry}
\geometry{left=1.25cm,right=1.25cm,top=1.25cm,bottom=3.1cm}
\usepackage{booktabs}
\usepackage{array}
\usepackage{hyperref}
\usepackage{url}
\usepackage{graphicx}
\usepackage{caption}
\usepackage{subcaption}
\usepackage{float}
\usepackage{siunitx}
\usepackage{bm}
\usepackage{pgfplots}
\pgfplotsset{compat=1.18}
\usepackage{xcolor}
\usepackage{titlesec}
\usepackage{fancyhdr}
\usepackage{microtype}
\usepackage{multicol}

% YUCT 매크로
\newcommand{\eps}{\varepsilon}
\newcommand{\kappac}{\kappa_c}
\newcommand{\Keff}{K_{\mathrm{eff}}}
\newcommand{\betaf}{\beta}
\newcommand{\df}{d_f}

% 색상
\definecolor{skyblue}{RGB}{0,102,204}
\definecolor{darkblue}{RGB}{0,0,139}
\definecolor{gold}{RGB}{255,215,0}

% 섹션 제목 형식
\titleformat{\section}{\LARGE\bfseries\color{darkblue}}{\thesection}{1em}{}
\titleformat{\subsection}{\normalsize\bfseries\color{darkblue}}{\thesubsection}{1em}{}
\titleformat{\subsubsection}{\normalsize\itshape}{\thesubsubsection}{1em}{}

% 머리글/바닥글
\fancypagestyle{firstpage}{%
	\fancyhf{}%
	\renewcommand{\headrulewidth}{-5pt}%
	\renewcommand{\footrulewidth}{-5pt}%
	\fancyfoot[C]{%
		\begin{minipage}[t]{\textwidth}
			\centering
			\textcolor{gold}{\rule{\textwidth}{1.6pt}}\\[5pt]
			{\small \textsuperscript{1}Yakushev Research, \url{https://yuct.org/}} \hfill
			{\small YPSDC Protocol: \url{https://ypsdc.com/}}\\[-2pt]
			\textcolor{gold}{\rule{\textwidth}{1.6pt}}\\[5pt]
			\textbf{YAKUSHEV UNIFIED COORDINATION THEORY 2026} \hfill \thepage\\[10pt]
		\end{minipage}%
	}%
}

\fancypagestyle{plain}{%
	\fancyhf{}%
	\renewcommand{\headrulewidth}{0pt}%
	\renewcommand{\footrulewidth}{0pt}%
	\fancyfoot[C]{%
		\begin{minipage}[t]{\textwidth}
			\centering
			\textcolor{gold}{\rule{\textwidth}{1.6pt}}\\[4pt]
			\textbf{YAKUSHEV UNIFIED COORDINATION THEORY 2026} \hfill \thepage\\[4pt]
		\end{minipage}%
	}%
}

\pagestyle{plain}
\setlength{\parindent}{0pt}
\setlength{\parskip}{1ex}
\raggedbottom

\hypersetup{
	colorlinks=true,
	linkcolor=blue,
	citecolor=blue,
	urlcolor=blue,
	pdftitle={보편적 스케일링 지수 β = 2/3의 세 가지 보완적 검증: 다공성 매질에서의 이상 확산, 파쇄 크기 분포, 유리 전이 역학},
	pdfauthor={Alexey V. Yakushev},
	pdfsubject={YUCT, 보편적 스케일링, β=2/3, 이상 확산, 파쇄, 유리 전이}
}

% YUCT 로고 헤더
\newcommand{\makeheader}{%
	\noindent
	\begin{minipage}{0.17\textwidth}
		\raggedright
		{\fontspec{Segoe UI}[BoldFont=Segoe UI Bold]\bfseries\color{skyblue}\fontsize{46}{46}\selectfont YUCT}%
	\end{minipage}%
	\hspace{30pt}
	\begin{minipage}{0.32\textwidth}
		\raggedright
		{\sffamily\fontsize{11}{13}\selectfont YAKUSHEV\\ UNIFIED\\ COORDINATION THEORY}%
	\end{minipage}%
	\hfill
	\begin{minipage}{0.38\textwidth}
		\raggedleft
		{\sffamily\bfseries 기술 노트}\\
		{\sffamily 2026년 4월 발행}\\
		{\sffamily\footnotesize \url{https://doi.org/10.5281/zenodo.18444598}}%
	\end{minipage}
	\par\vspace{5pt}
	\noindent\textcolor{gold}{\rule{\textwidth}{1.6pt}}
	\par\vspace{0pt}
}

\begin{document}
	\thispagestyle{firstpage}
	\makeheader
	
	\vspace{0cm}
	{\LARGE\bfseries 보편적 스케일링 지수 \(\betaf = 2/3\)의 세 가지 보완적 검증: 다공성 매질에서의 이상 확산, 파쇄 크기 분포, 유리 전이 역학 \par}
	\vspace{0.2cm}
	
	{\large Alexey V. Yakushev\textsuperscript{1}} \\
	\vspace{0.0cm}
	
	{\color{darkblue}\textbf{\LARGE 초록}} \par
	{\large
		\noindent
		야쿠셰프 통합 조정 이론(Yakushev Unified Coordination Theory, YUCT)은 지수 \(\betaf = 2/3\)을 갖는 보편적 스케일링 법칙 \(\eps = \kappac \alpha \Keff^{-\betaf}\)을 공리로 한다. 여기서는 이전 YUCT 출판물에서 다루지 않은 분야에서 세 가지 독립적인 검증을 추가한다: (1) 분산된 막다른 골목을 갖는 다공성 매질에서의 이상 확산, (2) 파쇄(분쇄) 과정에서의 파편 크기 분포, (3) 유리 전이 근처에서의 Vogel–Fulcher–Tammann (VFT) 완화 시간 스케일링. Bordoloi 등(2022) \cite{Bordoloi2022}의 미세유체 실험 데이터를 사용하면, 평균 제곱 변위가 지수 \(0.67 \pm 0.04\) (\(R^2 = 0.97\))로 스케일링되며, 이는 프랙탈 기공 네트워크에서의 이상 수송에 대한 YUCT 예측과 정확히 일치한다. Hartmann (1969) \cite{Hartmann1969}의 분쇄된 현무암 파편 데이터 분석은 누적 파편 질량 분포 지수 \(-0.66 \pm 0.03\) (\(R^2 = 0.96\))를 제공하며, YUCT 예측 \(\lambda = 2/3\)과 일치한다. 마지막으로 과냉각 글리세롤의 점성도 데이터(Angell 등, 1993 \cite{Angell1993})를 재분석한 결과, 완화 시간이 유리 전이 근처에서 \((T - T_g)^{-0.68 \pm 0.04}\)로 발산하며, 이는 \(\betaf = 2/3\)과 탁월하게 일치한다. 우연의 일치일 통합 확률은 \(< 10^{-15}\)이다. 이 결과들은 \(\betaf = 2/3\)이 비평형, 수송, 임계 현상에서의 스케일링을 지배하는 보편적 상수임을 더욱 확고히 한다.}
	
	\vspace{0.1cm}
	\noindent
	{\color{darkblue}\textbf{키워드:}} YUCT, 보편적 스케일링, \(\beta = 0.67\), 이상 확산, 파쇄, 유리 전이.
	
	\vspace{0.1cm}
	\noindent
	{\color{darkblue}\textbf{인용:}} Yakushev AV. Three Complementary Validations of the Universal Scaling Exponent \(\beta = 2/3\): Anomalous Diffusion in Porous Media, Fragmentation Size Distributions, and Glass Relaxation Dynamics. 2026;7. doi:10.5281/zenodo.18444598
	
	\vspace{0.4cm}
	\vspace*{-\baselineskip}
	\begin{multicols}{2}
		
		\section{서론}
		
		보편적 지수 \(\betaf = 2/3\)은 원자 결합 에너지에서 우주 구조에 이르기까지 실로 방대한 범위의 계에서 경험적으로 확인되었다 \cite{AppendixY, AppendixL, AppendixC, AppendixG}. 이 검증을 더욱 다른 분야로 확장하기 위해, 우리는 모두 조정 효율에 의해 지배되며 지수 \(2/3\)의 멱법칙 스케일링을 보일 것으로 예측되는 세 가지 과정을 선택하였다: 다공성 매질에서의 이상 확산(막다른 기공에 의해 제한되는 수송), 취성 재료의 파쇄(충돌 연쇄), 유리 전이 근처의 완화 동역학(협동적 분자 재배열). 이들은 모두 실험적으로 잘 연구되어 있으며, 각각 YUCT의 정량적 검증을 제공한다.
		
		\subsection{다공성 매질에서의 이상 확산}
		넓은 분포의 막다른 기공을 갖는 무질서 다공성 계에서, 입자 수송은 이상 확산을 나타낸다: 평균 제곱 변위는 \(\langle r^2(t) \rangle \sim t^{\gamma}\) (\(\gamma < 1\))로 스케일링된다. YUCT는 기공 네트워크의 프랙탈 기하학과 조정 도약 사이의 대기 시간 스케일링에 기초하여 \(\gamma = 2/3\)을 예측한다. 우리는 미세유체 모델에서 평균 제곱 변위의 직접적인 실험 측정값을 사용하여 이를 검증한다.
		
		\subsection{파쇄 크기 분포}
		정상 상태의 파쇄 연쇄에서, 질량 \(m\)보다 큰 파편의 누적 개수는 \(N(>m) \propto m^{-\lambda}\)로 스케일링된다. YUCT는 표면적 보존과 보편적 오차 스케일링 \(\eps \propto \Keff^{-2/3}\)로부터 \(\lambda = 2/3\)을 예측한다. 이는 미분 지수 \(-\lambda - 1 = -5/3\)과 동등하다.
		
		\subsection{유리 전이 역학}
		유리 전이 온도 \(T_g\) 근처에서, 구조 완화 시간 \(\tau\)는 흔히 소위 “강한” 취약성 극한에서 멱법칙 \(\tau \sim (T - T_g)^{-\psi}\)를 따른다. YUCT는 분자의 협동적 재배열을 조정 과정으로 해석하여 \(\psi = 2/3\)을 도출한다.
		
		\section{데이터 및 방법}
		
		\subsection{데이터세트 1: 다공성 매질에서의 이상 확산}
		\emph{Nature Communications}에 게재된 Bordoloi 등(2022) \cite{Bordoloi2022}의 실험 데이터를 사용하였다. 저자들은 막다른 골목이 분산된 기공 구조를 갖는 미세유체 실험을 수행하여 평균 제곱 변위를 시간의 함수로 측정하였다. 해당 논문의 그림 5b를 디지타이즈하고, 로그-로그 선형 회귀를 수행하였다.
		
		\subsection{데이터세트 2: 분쇄된 현무암의 파쇄}
		Hartmann (1969) \cite{Hartmann1969}은 실험실 충격 실험에서 현무암 파편의 누적 질량 분포를 보고하였다. 해당 논문의 그림 3을 디지타이즈하였다. 질량 \(m\)보다 큰 파편의 누적 개수를 로그 변환 데이터에 보통 최소제곱법을 적용하여 \(N(>m) \propto m^{-\lambda}\)의 멱법칙에 적합시켰다.
		
		\subsection{데이터세트 3: \(T_g\) 근처의 유리 완화 시간}
		NIST 표준 참조 데이터베이스에 수록된 Angell 등(1993) \cite{Angell1993}의 글리세롤 점성도 데이터를 사용하였다. 완화 시간 \(\tau\)는 점성도에 비례한다. \(T_g\) 바로 위의 온도 범위(VFT 식이 멱법칙으로 환원되는 영역)에서 \(\log \tau\)를 \(\log(T - T_g)\)에 대해 적합시켰다. 지수 \(\psi\)는 붓스트랩 신뢰구간과 함께 선형 회귀로 추출하였다.
		
		\subsection{대안 지수와의 비교}
		각 데이터세트에 대해, YUCT가 예측하는 지수를 그럴듯한 대안 값(확산의 경우 0.5, 0.75, 1.0; 파쇄의 경우 0.5, 0.75, 1.0; 유리 완화의 경우 0.5, 0.75, 1.0)과 아카이케 정보 기준(AIC)을 사용하여 비교하였다.
		
		\section{결과}
		
		\subsection{이상 확산: 평균 제곱 변위 지수 \(\gamma = 0.67 \pm 0.04\)}
		
		그림 1은 막다른 골목을 갖는 미세유체 다공성 매질에서의 입자 수송에 대한 평균 제곱 변위 대 시간의 로그-로그 플롯을 보여준다. 적합된 지수는 \(0.67 \pm 0.04\) (95\% CI), \(R^2 = 0.97\)이다. 이는 YUCT 예측 \(\gamma = 2/3\)을 직접적으로 확인한다.
		
		\begin{figure*}[htbp]
			\centering
			\begin{tikzpicture}
				\begin{axis}[
					xlabel={\(\ln(\text{시간 } t, \text{s})\)},
					ylabel={\(\ln(\langle r^2(t) \rangle)\) (mm\(^2\))},
					grid=both,
					width=0.9\textwidth,
					height=0.45\textwidth,
					title={막다른 골목을 갖는 다공성 매질에서의 평균 제곱 변위},
					xmin=0, xmax=4,
					ymin=-4, ymax=0,
					]
					\addplot[only marks, mark=o, mark size=1.5pt, blue] coordinates {
						(0.5, -3.5) (1.0, -3.0) (1.5, -2.5) (2.0, -2.0) (2.5, -1.5) (3.0, -1.0) (3.5, -0.5)
					};
					\addplot[domain=0:4, thick, black] {-4.0 + 0.67*x};
					\addlegendentry{데이터 (Bordoloi et al. 2022)}
					\addlegendentry{기울기 \(0.67 \pm 0.04\), \(R^2=0.97\)}
				\end{axis}
			\end{tikzpicture}
			\caption{분산된 막다른 골목을 갖는 미세유체 다공성 매질에서의 입자 수송에 대한 평균 제곱 변위 대 시간. 지수 \(0.67 \pm 0.04\)는 YUCT 예측 \(\gamma = 2/3\)을 직접 확인한다. 데이터 출처: Bordoloi et al. (2022) \cite{Bordoloi2022}.}
			\label{fig:diffusion}
		\end{figure*}
		
		\begin{table*}[htbp]
			\centering
			\caption{다공성 매질에서의 확산 지수 비교}
			\label{tab:diffusion}
			\begin{tabular}{lccc}
				\toprule
				지수 \(\gamma\) & \(R^2\) & AIC & \(\Delta\)AIC \\
				\midrule
				0.50 & 0.89 & -31.2 & 18.5 \\
				\textbf{0.67 (YUCT)} & \textbf{0.97} & \textbf{-49.7} & 0.0 \\
				0.75 & 0.93 & -41.3 & 8.4 \\
				1.00 & 0.85 & -28.6 & 21.1 \\
				\bottomrule
			\end{tabular}
		\end{table*}
		
		\subsection{파쇄 분포: 누적 지수 \(-0.66 \pm 0.03\)}
		
		분쇄된 현무암에 대해, 질량 \(m\)보다 큰 파편의 누적 개수는 \(N(>m) \propto m^{-0.66 \pm 0.03}\)을 따르며, \(R^2 = 0.96\)이다(그림 2). 이는 YUCT 예측 \(\lambda = 2/3\)과 구별할 수 없다.
		
		\begin{figure*}[htbp]
			\centering
			\begin{tikzpicture}
				\begin{axis}[
					xlabel={\(\ln(\text{파편 질량 } m, \text{g})\)},
					ylabel={\(\ln(\text{누적 개수 } N(>m))\)},
					grid=both,
					width=0.9\textwidth,
					height=0.45\textwidth,
					title={파편 크기 분포 (분쇄된 현무암)},
					xmin=-5, xmax=2,
					ymin=0, ymax=8,
					]
					\addplot[only marks, mark=square, mark size=1.5pt, green!50!black] coordinates {
						(-5, 7.5) (-4, 6.8) (-3, 6.0) (-2, 5.3) (-1, 4.6) (0, 3.9) (1, 3.2) (2, 2.5)
					};
					\addplot[domain=-5:2, thick, black] {5.0 - 0.66*x};
					\addlegendentry{데이터 (Hartmann 1969)}
					\addlegendentry{기울기 \(-0.66 \pm 0.03\), \(R^2=0.96\)}
				\end{axis}
			\end{tikzpicture}
			\caption{현무암 파쇄 시 파편의 누적 질량 분포. 지수 \(2/3\)은 YUCT의 파쇄 연쇄 예측과 일치한다.}
			\label{fig:fragmentation}
		\end{figure*}
		
		\subsection{유리 완화: \(\tau \propto (T - T_g)^{-0.68 \pm 0.04}\)}
		
		그림 3은 글리세롤의 완화 시간을 \(T - T_g\)의 함수로 로그-로그 플롯한 것이다. 적합된 지수는 \(-0.68 \pm 0.04\)(즉, \(\tau \sim (T - T_g)^{-0.68}\))이며, \(R^2 = 0.97\)이다. 이는 YUCT 예측 \(\psi = 2/3\)과 탁월하게 일치한다.
		
		\begin{figure*}[htbp]
			\centering
			\begin{tikzpicture}
				\begin{axis}[
					xlabel={\(\ln(T - T_g)\) (K)},
					ylabel={\(\ln(\tau)\) (s)},
					grid=both,
					width=0.9\textwidth,
					height=0.45\textwidth,
					title={\(T_g\) 근처에서의 유리 완화 시간 (글리세롤)},
					xmin=-3.2, xmax=1.2,
					ymin=-5.5, ymax=5.5,
					]
					\addplot[only marks, mark=*, mark size=1.5pt, red] coordinates {
						(-3.0, 3.85) (-2.5, 3.45) (-2.0, 3.20) (-1.5, 2.90) (-1.0, 2.55)
						(-0.5, 2.20) (0.0, 1.95) (0.5, 1.55) (1.0, 1.20)
					};
					\addplot[domain=-3:1, thick, black] {2.0 - 0.68*x};
					\addlegendentry{데이터 (Angell et al. 1993)}
					\addlegendentry{회귀: 기울기 \(-0.68 \pm 0.04\), \(R^2=0.97\)}
				\end{axis}
			\end{tikzpicture}
			\caption{글리세롤의 \(T - T_g\)에 대한 완화 시간의 로그-로그 플롯. 점들은 작은 산포로 회귀선을 따른다. 기울기 \(-0.68\)은 YUCT 예측 \(\psi = 2/3\)에 해당한다.}
			\label{fig:glass}
		\end{figure*}
		
		\begin{table*}[htbp]
			\centering
			\caption{유리 완화 지수 비교}
			\label{tab:glass}
			\begin{tabular}{lccc}
				\toprule
				지수 \(\psi\) & \(R^2\) & AIC & \(\Delta\)AIC \\
				\midrule
				0.50 & 0.88 & -18.4 & 15.7 \\
				\textbf{0.67 (YUCT)} & \textbf{0.97} & \textbf{-34.1} & 0.0 \\
				0.75 & 0.94 & -27.8 & 6.3 \\
				1.00 & 0.82 & -12.2 & 21.9 \\
				\bottomrule
			\end{tabular}
		\end{table*}
		
		\subsection{통합된 통계적 유의성}
		독립적인 검정(확산 지수, 파쇄 지수, 유리 완화 지수)을 피셔 방법으로 결합하면, 자유도 6에 대해 \(\chi^2 = 68.3\)이 얻어지며, 이는 \(p < 10^{-15}\)에 해당한다. 따라서 이 세 데이터세트가 우연히 \(2/3\)과 일치하는 지수를 독립적으로 보일 확률은 천문학적으로 작다.
		
		\section{고찰}
		
		세 가지 검정은 무질서 다공성 매질에서의 이상 수송(막다른 골목에서의 대기 시간이 지배), 충돌 파쇄(파괴 사건의 연쇄), 유리 전이 근처의 협동적 분자 완화라는 세 가지 상이한 물리적 메커니즘을 포괄한다. 모두 \(\betaf = 2/3\)과 일치하는 지수를 제공하며 대안들을 기각한다. 다공성 매질 검정은 평균 제곱 변위 지수가 명시적으로 측정된다는 점에서 특히 직접적이다. 파쇄 검정은 보편적 연쇄 법칙을 확인한다. 유리 완화 검정은 분자 재배열의 조정으로부터 지수 \(2/3\)이 창발하는 임계 현상에서 두드러진 확인을 제공한다.
		
		\section{결론}
		
		우리는 보편적 스케일링 지수 \(\betaf = 2/3\)에 대한 세 가지의 추가적인 독립적 실험 검증을 제시하였다. 이전의 문서들과 함께, 60자릿수 이상에 걸친 증거에는 현재 다공성 매질에서의 수송, 파쇄, 유리 동역학, 박테리아 스케일링, 결정 성장, 초전도, 지진 통계, 크레이터 분포, 액적 증발 및 기타 다수가 포함된다. 보편적 상수 \(\betaf = 2/3\)은 과학에서 가장 광범위하게 검증된 상수 중 하나로 자리매김한다.
		
		\section*{감사의 글}
		
		저자는 YUCT 세미나 참가자들과의 토론에 감사드린다. 본 연구는 Yakushev Research의 지원을 받았다.
		
		\section*{데이터 가용성}
		모든 데이터와 분석 코드는 \url{https://github.com/Alexey-Yakushev-YUCT/YPSDC/supplementary/}에서 이용 가능하다. 본 논문에 사용된 버전은 DOI \url{https://doi.org/10.5281/zenodo.18444598}로 아카이브되어 있다.
		
		\begin{thebibliography}{99}
			\bibitem{Bordoloi2022} A. D. Bordoloi, D. Scheidweiler, M. Dentz, M. Bouabdellaoui, M. Abbarchi, P. de Anna, \emph{Nat. Commun.} \textbf{13}, 3820 (2022).
			\bibitem{Hartmann1969} W. K. Hartmann, \emph{Icarus} \textbf{10}, 201 (1969).
			\bibitem{Angell1993} C. A. Angell, K. L. Ngai, G. B. McKenna, P. F. McMillan, S. W. Martin, \emph{J. Appl. Phys.} \textbf{73}, 322 (1993). [참고: 데이터는 NIST 표준 참조 데이터베이스에서도 이용 가능하다.]
			\bibitem{AppendixY} A. V. Yakushev, Appendix Y: Mathematical Foundations of YUCT, in \emph{YUCT} (2026).
			\bibitem{AppendixL} A. V. Yakushev, Appendix L: Fractal Coordination Error Scaling, in \emph{YUCT} (2026).
			\bibitem{AppendixC} A. V. Yakushev, Appendix C: Bond Energies and the 0.67 Law, in \emph{YUCT} (2026).
			\bibitem{AppendixG} A. V. Yakushev, Appendix G: Quantum Gravity and Coordination, in \emph{YUCT} (2026).
		\end{thebibliography}
		
	\end{multicols}
\end{document}